\chapter{Protocol Layers}

\section{Why Three Protocols?}

Agent communication in production spans three distinct contexts, each with
different requirements:

\begin{itemize}[leftmargin=1.5em]
  \item \textbf{Native intra-node} — Agents on the same Connector node need
        fast, capability-gated, typed messaging with full auditability.
        This is what \textbf{CNP} (Connector Native Protocol) is designed for.

  \item \textbf{Cross-platform} — An agent may need to communicate with an agent
        running on a different platform (LangChain, AutoGen, another Connector node).
        The \textbf{A2A} (Agent-to-Agent) bridge handles cross-platform interop.

  \item \textbf{External tools} — Agents need access to databases, APIs, search
        engines, and other services. The \textbf{MCP} (Model Context Protocol)
        bridge connects to any MCP-compatible tool server.
\end{itemize}

\begin{figure}[H]
\centering
\begin{tikzpicture}[
  box/.style={rectangle, rounded corners=4pt, draw=#1!70, fill=#1!9,
              text width=3.2cm, align=center, minimum height=1.3cm,
              font=\small\bfseries, inner sep=5pt},
  bridge/.style={rectangle, rounded corners=3pt, draw=gray!50, fill=gray!6,
                 text width=2.4cm, align=center, minimum height=0.9cm,
                 font=\footnotesize\bfseries},
  node distance=1.5cm and 2cm
]
\node[box=CKernel] (agent) {Connector\\Agent};

\node[box=CProto, above right=1cm and 2.5cm of agent] (cnp)
  {CNP\\{\footnotesize Native 7-layer}};
\node[box=CDX, right=2.5cm of agent] (a2a)
  {A2A Bridge\\{\footnotesize Cross-platform}};
\node[box=CBiz, below right=1cm and 2.5cm of agent] (mcp)
  {MCP Bridge\\{\footnotesize External tools}};

\node[box=CExec, right=2.2cm of cnp] (peer)   {Peer\\Agent};
\node[box=CExec, right=2.2cm of a2a] (extern) {External\\Agent};
\node[box=CExec, right=2.2cm of mcp] (tool)   {Tool Server};

\draw[arr=CProto] (agent) -- (cnp);
\draw[arr=CDX]    (agent) -- (a2a);
\draw[arr=CBiz]   (agent) -- (mcp);
\draw[arr=CProto] (cnp) -- (peer);
\draw[arr=CDX]    (a2a) -- (extern);
\draw[arr=CBiz]   (mcp) -- (tool);

\node[above=0.1cm of cnp, font=\scriptsize\itshape, text=CProto]
  {typed, capability-gated};
\node[above=0.1cm of a2a, font=\scriptsize\itshape, text=CDX]
  {DID-addressed};
\node[below=0.1cm of mcp, font=\scriptsize\itshape, text=CBiz]
  {JSON-RPC, stdio};

\end{tikzpicture}
\caption{Connector protocol landscape — CNP (native), A2A (cross-platform), MCP (tools)}
\label{fig:protocol-landscape}
\end{figure}

\section{CNP: Connector Native Protocol}

CNP is the primary communication protocol for agents running on the same
Connector node. It is modelled on the OSI 7-layer model, but every layer is
designed specifically for agent-to-agent communication: content addressing,
capability gating, and cryptographic receipts are first-class concerns.

\subsection{The 7-Layer Stack}

\begin{figure}[H]
\centering
\begin{tikzpicture}[node distance=0.45cm]
  \node[layerbox=CBiz]
    {L7 — Application\quad Message types, payload routing, application semantics};
  \node[layerbox=CDX, below=0.45cm of current bounding box.south]
    {L6 — Session\quad Session establishment, lifecycle, teardown, receipts};
  \node[layerbox=CDX, below=0.45cm of current bounding box.south]
    {L5 — Routing\quad Cross-cell routing, delivery guarantees, dead-letter queue};
  \node[layerbox=CLogic, below=0.45cm of current bounding box.south]
    {L4 — Port Management\quad Typed channels, capability gating, buffer management};
  \node[layerbox=CProto, below=0.45cm of current bounding box.south]
    {L3 — Security\quad HMAC signatures, anti-replay tokens, TTL enforcement};
  \node[layerbox=CProto, below=0.45cm of current bounding box.south]
    {L2 — Encryption\quad AEAD encryption (ChaCha20-Poly1305), nonce management};
  \node[layerbox=CKernel, below=0.45cm of current bounding box.south]
    {L1 — Content-Addressed Codec\quad CID computation, size validation, serialisation};
\end{tikzpicture}
\caption{CNP 7-layer stack — modelled on OSI, purpose-built for agents}
\label{fig:cnp-stack}
\end{figure}

Every message passes through all seven layers on the way out, and all seven
in reverse on the way in. Layer 1 gives tamper detection; Layer 2 gives
confidentiality; Layer 3 gives authentication and replay prevention; Layer 4
gives capability enforcement; Layers 5--7 give routing, sessions, and
application semantics.

\subsection{Port-Based Communication}

CNP does not use raw sockets or queues. All communication happens through
\textbf{typed Ports} — named, capability-gated channels attached to an agent.

\begin{table}[H]
\centering
\renewcommand{\arraystretch}{1.4}
\begin{tabularx}{\textwidth}{@{}llX@{}}
\toprule
\textbf{Port Type} & \textbf{Direction} & \textbf{Used for} \\
\midrule
\texttt{Sensor}    & Inbound  & Receiving raw sensor/environmental data \\
\texttt{Actuator}  & Outbound & Sending commands to physical or virtual actuators \\
\texttt{Memory}    & Both     & Sharing MemPackets between agents \\
\texttt{Tool}      & Outbound & Invoking tools and receiving results \\
\texttt{Cognitive} & Both     & Passing cognitive messages between agents \\
\texttt{Generic}   & Both     & Untyped fallback for custom payloads \\
\bottomrule
\end{tabularx}
\caption{CNP port types and their communication semantics}
\end{table}

Opening a port requires a \textbf{Capability} token signed by the issuing agent.
No capability, no port — the kernel rejects unauthorised connection attempts at
Layer 4 before any payload is processed.

\subsection{Capability Attenuation}

Capabilities can be \emph{delegated} (shared with another agent) but only with
equal or reduced scope — never expanded. This UCAN-style attenuation forms a
trust chain:

\begin{figure}[H]
\centering
\begin{tikzpicture}[
  capnode/.style={rectangle, rounded corners=3pt, draw=CLogic!70, fill=CLogic!8,
              text width=3.8cm, align=left, minimum height=1.4cm,
              font=\footnotesize, inner sep=6pt},
  node distance=1.8cm
]
\node[capnode] (root) {
  \textbf{Root capability}\\
  issuer: platform\\
  actions: [read, write, exec]\\
  depth: 0
};
\node[capnode, right=of root] (d1) {
  \textbf{Delegated (depth 1)}\\
  issuer: agent-A\\
  actions: [read, write]\\
  depth: 1
};
\node[capnode, right=of d1] (d2) {
  \textbf{Delegated (depth 2)}\\
  issuer: agent-B\\
  actions: [read]\\
  depth: 2
};
\node[rectangle, draw=CDanger!60, fill=CDanger!6, rounded corners=3pt,
      text width=2.2cm, align=center, minimum height=1.0cm, font=\footnotesize,
      right=of d2] (reject) {\textbf{Depth 3}\\rejected\\by kernel};

\draw[arr=CLogic] (root) -- node[above, font=\scriptsize]{\textit{attenuate}} (d1);
\draw[arr=CLogic] (d1)   -- node[above, font=\scriptsize]{\textit{attenuate}} (d2);
\draw[arr=CDanger, dashed] (d2) -- (reject);

\end{tikzpicture}
\caption{Capability attenuation chain — scope can only be reduced, never expanded; max depth is 3}
\label{fig:capability-chain}
\end{figure}

\subsection{Delivery Guarantees}

CNP supports three delivery modes, selectable per message:

\begin{table}[H]
\centering
\renewcommand{\arraystretch}{1.45}
\begin{tabularx}{\textwidth}{@{}llXl@{}}
\toprule
\textbf{Mode} & \textbf{Retries} & \textbf{Behaviour} & \textbf{Use case} \\
\midrule
\texttt{AtMostOnce}   & 0 & Fire and forget; no acknowledgement required
  & Telemetry, non-critical events \\
\texttt{AtLeastOnce}  & $\infty$ & Retry until acked; possible duplicates
  & Commands requiring eventual delivery \\
\texttt{ExactlyOnce}  & $\infty$ & Idempotent delivery with deduplication key
  & Financial transactions, state mutations \\
\bottomrule
\end{tabularx}
\caption{CNP delivery guarantee modes}
\end{table}

Undeliverable messages go to a per-agent \textbf{dead-letter queue}, where they
can be inspected and replayed by an operator via \texttt{connectorctl}.

\section{A2A: Agent-to-Agent Bridge}

The A2A bridge allows Connector agents to communicate with agents on external
platforms — LangChain, AutoGen, CrewAI, or other Connector nodes — without
requiring CNP on both sides.

The bridge is \textbf{bidirectional}: Connector can receive inbound A2A messages
and convert them to CNP for internal routing, and it can convert outbound CNP
messages to A2A for delivery to external agents.

\begin{figure}[H]
\centering
\begin{tikzpicture}[
  box/.style={rectangle, rounded corners=4pt, draw=#1!70, fill=#1!9,
              text width=3cm, align=center, minimum height=1.1cm,
              font=\small\bfseries},
  bridge/.style={rectangle, rounded corners=3pt, draw=CDX!60, fill=CDX!8,
                 text width=3.2cm, align=center, minimum height=1.1cm,
                 font=\small\bfseries},
  node distance=1.2cm
]
\node[box=CKernel] (cn) {Connector\\Agent (CNP)};
\node[bridge, right=2cm of cn] (br) {A2A $\Longleftrightarrow$ CNP\\Bridge};
\node[box=CExec, right=2cm of br] (ext) {External\\Agent (A2A)};

\draw[darr=CDX] (cn) -- (br);
\draw[darr=CDX] (br) -- (ext);

\node[below=0.2cm of br, font=\footnotesize\itshape, text=gray]
  {DID translation, payload mapping};

\end{tikzpicture}
\caption{A2A bridge — bidirectional protocol translation between CNP and A2A}
\end{figure}

All inbound A2A messages are logged as MemPackets by the bridge, so even
cross-platform communication is auditable from the Connector side.

\section{MCP: Model Context Protocol Bridge}

MCP is an open standard for connecting AI systems to external tool and data
providers. Connector implements an MCP client bridge that allows any registered
MCP server to be used as a tool by Connector agents.

\begin{figure}[H]
\centering
\begin{tikzpicture}[
  box/.style={rectangle, rounded corners=4pt, draw=#1!70, fill=#1!9,
              text width=3cm, align=center, minimum height=1.1cm,
              font=\small\bfseries},
  node distance=1.6cm
]
\node[box=CKernel] (ag) {Agent\\(CLS contract)};
\node[box=CBiz, right=of ag] (bridge) {MCP Bridge\\(L4 tool gate)};
\node[box=CExec, right=of bridge] (server) {MCP Server\\(stdio/HTTP)};
\node[box=gray, right=of server] (tool) {Database\\API / Search};

\draw[arr=CKernel!60] (ag)     -- node[above, font=\scriptsize]{\texttt{tool\_call}} (bridge);
\draw[arr=CBiz!60]    (bridge) -- node[above, font=\scriptsize]{JSON-RPC} (server);
\draw[arr=CExec!60]   (server) -- (tool);
\draw[arr=CExec!60, bend left=20] (tool.north) to (server.north);
\draw[arr=CBiz!60, bend left=20]  (server.north) to (bridge.north);
\draw[arr=CKernel!60, bend left=20] (bridge.north) to
  node[above, font=\scriptsize]{MemPacket} (ag.north);

\end{tikzpicture}
\caption{MCP tool invocation flow — every call is wrapped in a MemPacket for auditability}
\end{figure}

MCP servers are configured at node startup. Each registered server exposes
a list of named tools with JSON Schema input specifications. Agents address
tools by name; the bridge handles the protocol translation transparently.

\section{Protocol Selection Guide}

\begin{table}[H]
\centering
\renewcommand{\arraystretch}{1.45}
\rowcolors{2}{gray!5}{white}
\begin{tabularx}{\textwidth}{@{}Xccc@{}}
\toprule
\textbf{Use case} &
  \textbf{CNP} & \textbf{A2A} & \textbf{MCP} \\
\midrule
Intra-node agent-to-agent messaging   & \checkmark & & \\
Cross-platform agent communication    & & \checkmark & \\
External tool / API integration       & & & \checkmark \\
Capability-gated typed channels       & \checkmark & & \\
Typed sensor / actuator data          & \checkmark & & \\
Cognitive reasoning message passing   & \checkmark & \checkmark & \\
MemPacket sharing between agents      & \checkmark & \checkmark & \\
Memory knowledge base queries         & \checkmark & & \checkmark \\
Cross-node cluster communication      & \checkmark & \checkmark & \\
\bottomrule
\end{tabularx}
\caption{Protocol selection guide — choose based on topology and capability requirements}
\end{table}

\section{Shared Security Properties}

Regardless of which protocol is used, all three provide a common set of
security guarantees enforced at the kernel level:

\begin{table}[H]
\centering
\renewcommand{\arraystretch}{1.4}
\begin{tabularx}{\textwidth}{@{}lX@{}}
\toprule
\textbf{Property} & \textbf{How it is enforced} \\
\midrule
\textbf{Authentication}    & Sender DID verified via Ed25519 signature on every message \\
\textbf{Integrity}         & L1 CID detects any in-flight tampering; L3 HMAC prevents forgery \\
\textbf{Confidentiality}   & L2 AEAD encryption (CNP); TLS for A2A and MCP \\
\textbf{Non-repudiation}   & Signed delivery receipts stored as MemPackets \\
\textbf{Audit trail}       & Every protocol operation logged to the Memory Kernel \\
\bottomrule
\end{tabularx}
\caption{Security properties shared across all three protocol systems}
\end{table}
